%=====================================================================
\section{Generator plug-ins (v2.0)}\label{sec:v20gen}
% numbered from 700 so that the parallel 2.x threads (106-108 for the IR, 600-605 for privacy) do not collide
\newcounter{savedpropgen}\setcounter{savedpropgen}{\value{proposition}}
\setcounter{proposition}{699}
Version 2.0 separates the \emph{relational skeleton} of a synthetic database from the \emph{attribute columns} of
each table (\texttt{relsynth/generators}, guide \texttt{docs/generators.md}, design record
\texttt{docs/adr/0002-generator-plugins.md}). The skeleton $\mathcal S$---which rows exist, their keys, the links
between them, the latent segment of every row, counts and scale---is drawn by the engine of
Sections~\ref{sec:star}--\ref{sec:v10}. The attribute columns $x_t$ of table $t$ are drawn by a \emph{generator}
$G_t$ chosen in the spec: a pair of maps $\mathrm{fit}$ and $\mathrm{sample}$, where $\mathrm{sample}$ returns, for
rows with contexts $c_i$, values $x_i\sim q_{G_t}(\cdot\mid c_i;\theta_t)$. The context $c_i=(g_i,\bar g_i,a_i)$ is
the row's latent group $g_i$ (its segment, or its segment block for a fact), the coarser group $\bar g_i$ the engine
backs off to, and---for a generator that declares \emph{parent} conditioning---the attributes $a_i$ of the row's
synthetic parents. Write $\mathcal R$ for the rules of the spec, split into \emph{computed} rules (derive,
aggregate, cumulative, computed columns) and \emph{check} rules and unique attribute keys. The checks are in
\texttt{tests/unit/test\_generators.py}, and the identity runs and scorecards of every generator on every suite are
in the T03 report.

\subsection{What the engine guarantees with any generator}
\begin{proposition}[Skeleton and rules do not depend on the generator]\label{prop:genskeleton}
Let every $G_t$ satisfy the contract: $\mathrm{sample}(n,c,\rho)$ returns $n$ rows with exactly the table's own
attribute columns. Then, whatever the generators,
\begin{enumerate}
\item every synthetic dataset has unique, non-null primary keys; every foreign key is null or names an existing
parent row; trees are acyclic; and the row counts, scaling and link structure follow the laws of
Sections~\ref{sec:star}--\ref{sec:v10};
\item every computed rule holds exactly on every row;
\item a row that breaks a check rule or a unique attribute key after generation is redrawn from
$q_{G_t}(\cdot\mid c_i;\theta_t)$ until it holds or the round limit is reached. If the acceptance probability
$\pi_i=q_{G_t}(\mathcal R_i\mid c_i)$ is positive, each accepted row has law $q_{G_t}(\cdot\mid c_i,\mathcal R_i)$,
and the probability that a row still breaks a rule after $r$ rounds is $(1-\pi_i)^{r+1}$; such rows are reported.
\end{enumerate}
Parts of the skeleton that read attributes (rule-constrained links, key order by an attribute, attribute strata,
group effects, type 2 history) are functions of the generated values. Their law is then the 1.0 law given those
values.
\end{proposition}
\begin{proof}
(1) Keys and links are produced by the key and link samplers (Props.~\ref{prop:eta}--\ref{prop:cumulative} and the
v0.4--v1.0 sections). They read the fitted relational model and the entity encodings used for segment discovery,
which are built before and independently of any generator (\texttt{TableModel} encodes the columns; the generator is
attached afterwards). The generator's output is written into attribute columns only. The engine checks that it
returns exactly those columns (\texttt{TableModel.sample}), so no key or link column can be overwritten. (2) Computed
columns are evaluated by the engine after generation in the topological order of their inputs (Prop. on computed
columns, v0.7), from the final values, so they hold whatever the inputs. (3) The rejection loop
(\texttt{Synthesizer.\_finish}) redraws the offending rows by calling the same generator with the same context. The
rounds are independent draws from $q(\cdot\mid c_i)$. By the standard rejection-sampling argument (von Neumann
1951), the first accepted draw has the restricted law, and the number of failed rounds is geometric with success
probability $\pi_i$. The last sentence holds because these blocks are evaluated by the 1.0 code on the generated
columns.
\end{proof}

\begin{proposition}[Seed determinism]\label{prop:genseed}
If each generator draws every random number from the generator $\rho$ it is handed (or from generators seeded from
$\rho$), and its fit and sample are deterministic functions of their inputs and the state of $\rho$, then the fitted
model and every synthetic dataset are deterministic functions of the normalised spec, the data and the seed.
\end{proposition}
\begin{proof}
The engine's fitting stream is seeded by the spec seed, and the stream of dataset $m$ by (seed, sample seed, $m$)
(\texttt{posterior.run\_streams}). It hands $\rho$ to every generator call in a fixed order: tables in spec order,
facts in fact order, rejection rounds in row order. A composition of deterministic maps of
$(\text{inputs},\rho\text{-state})$ is deterministic, and so is the state of $\rho$ after each call, by induction on
the calls. The hypothesis is the plug-in's responsibility. The built-in generators satisfy it (numpy only, no global
state), and on a GPU it requires deterministic kernels.
\end{proof}

\begin{proposition}[The copula plug-in is relsynth 1.0]\label{prop:gencopula}
With every table on the \texttt{copula} generator (the default), the fitted model and every synthetic dataset are
identical, byte for byte, to relsynth 1.0's for the same spec, data and seed.
\end{proposition}
\begin{proof}
The copula generator consists of 1.0's \texttt{TableModel} methods \texttt{fit}, \texttt{\_fit\_marginals},
\texttt{redraw} and \texttt{sample}, moved without changing any operation. It reads the same encoded arrays (the
encoding stays in \texttt{TableModel} and runs before the generator is attached). It calls the same random
functions on the same $\rho$ in the same order. The order of the per-dataset redraw is unchanged: marginals, then
copula matrices, then the weights of pooled tail levels. The decoding loop is the same code. The engine adds no
random call when it attaches a generator, and the registry and compute resolution draw no random numbers. Hence
every intermediate value, and the output, is the same. The identity was also checked by running unmodified
\texttt{main} and this version on the eleven fake scenarios (two posterior datasets each), on Chinook (in memory, and
pushed down with two datasets), on Contoso (SQLite and Parquet), on bi/COREP and on AdventureWorks DW, and comparing
the SHA-256 of every output file.
\end{proof}

\begin{proposition}[Posterior predictive datasets]\label{prop:genpost}
Suppose that each generator's \texttt{posterior\_draw} replaces $\theta_t$ by an independent draw from
$p(\theta_t\mid x_t,\hat{\mathcal S})$, its posterior given the table's data and the fitted partition $\hat{\mathcal
S}$. Then dataset $m$ is a draw from the posterior predictive distribution of the joint model (skeleton and
attribute blocks) given the fitted partition, as in 1.0 (Prop.~\ref{prop:ppc}). If
\texttt{posterior\_draw} draws from an approximation $\tilde p(\theta_t)$, the dataset is a draw from the
corresponding approximate predictive. If it does nothing (\texttt{posterior=False}), datasets share $\theta_t$ and
differ only by sampling noise and by the relational parameters.
\end{proposition}
\begin{proof}
For each dataset, the engine draws the relational parameters from their posterior given $\hat{\mathcal S}$ and then
calls \texttt{posterior\_draw} on every fitted table with the dataset's parameter stream
(\texttt{Synthesizer.redraw}). The attribute blocks are conditionally independent of the relational parameters
given the partition and the contexts: the joint posterior factorises over the blocks, because each block's
likelihood involves only its own parameters and the contexts. A joint draw is therefore obtained by drawing each
block independently. Rows are then sampled given the drawn parameters, which is the definition of a posterior
predictive draw (Rubin 1993). The two weaker cases replace the block's exact posterior by $\tilde p$, or by a point
mass at the fitted value.
\end{proof}
\noindent\emph{The built-in generators.}
\begin{itemize}
\item \texttt{copula} is exact in this sense (1.0).
\item \texttt{independent} uses the copula's marginal posteriors and is exact for its own model.
\item \texttt{cart} redraws Bayesian-bootstrap leaf weights given trees fitted once. That is the exact posterior of
the leaf distributions under a Dirichlet process with vanishing concentration given the tree (Rubin 1981, Prop.~\ref{prop:cartbb}),
and an approximation of the posterior over trees. synthpop's ``proper'' synthesis refits the trees on a bootstrap
sample instead.
\end{itemize}

\subsection{Conditioning}
\begin{proposition}[What group-only conditioning keeps and loses]\label{prop:gencond}
Let a fact row $i$ with group $g_i$ have parent attributes $a_i$ and own attributes $x_i$. If the generator
conditions on the group only, then in the synthetic data $x_i\perp a_i\mid g_i$, and
\[\operatorname{Cov}(x_i,a_i)=\operatorname{Cov}_{g}\bigl(\mathbb E[x\mid g],\mathbb E[a\mid g]\bigr),\]
the covariance between the group means. With parent conditioning, $q(x\mid g,a)$ can also carry the within-group
dependence on $a$.
\end{proposition}
\begin{proof}
The engine draws the skeleton (and with it $g_i$ and the parent rows, hence $a_i$) first, then $x_i$ from
$q(\cdot\mid g_i)$, which does not read $a_i$; so $x_i$ is conditionally independent of $a_i$ given $g_i$. The law of
total covariance gives $\operatorname{Cov}(x,a)=\mathbb E_g\operatorname{Cov}(x,a\mid g)+\operatorname{Cov}_g(\mathbb
E[x\mid g],\mathbb E[a\mid g])$, and the first term is zero by conditional independence.
\end{proof}

\begin{proposition}[Independent marginals are a latent class model]\label{prop:genind}
The \texttt{independent} generator has the copula generator's marginal posteriors per group and the identity
correlation. Within a group its columns are independent, and between columns $j,k$ of a table
$\operatorname{Cov}(x_j,x_k)=\operatorname{Cov}_g(\mathbb E[x_j\mid g],\mathbb E[x_k\mid g])$. It is therefore the
latent class model of Goodman (1974), with the latent segments as classes and the hierarchical Dirichlet marginals of
Section~\ref{sec:star} (Bayesian form: Dunson and Xing 2009).
\end{proposition}
\begin{proof}
It fits the marginals with the copula generator's code (\texttt{SegmentMarginals}). It draws each column as
$F^{-1}_{g,j}(U_j)$ with independent uniforms $U_j$. That is the copula sampler with $C=I$, because $\Phi(Z_j)$ with
$Z\sim N(0,I)$ are independent uniforms. Independence within a group follows, and the covariance identity is
Prop.~\ref{prop:gencond} applied to two columns of the same row.
\end{proof}

\subsection{Sequential CART}
\begin{proposition}[Leaf donors with Bayesian-bootstrap weights]\label{prop:cartbb}
Let leaf $\ell$ of the tree of column $j$ hold the observed values $y_1,\dots,y_{n_\ell}$, with
$n_\ell\ge$\,\texttt{min\_leaf}. Draw $w\sim\mathrm{Dirichlet}(1,\dots,1)$ once per dataset and draw a synthetic
value $Y=y_I$ with $P(I=i\mid w)=w_i$. Then:
\begin{enumerate}
\item marginally over $w$, $Y$ has the empirical law of the leaf, $P(Y=v)=\#\{i:y_i=v\}/n_\ell$;
\item given $w$, the probability of a value $v$ with count $c$ in the leaf has mean $c/n_\ell$ and variance
$c(n_\ell-c)/\bigl(n_\ell^2(n_\ell+1)\bigr)$;
\item $w$ is the posterior of the leaf's distribution under a Dirichlet process prior whose concentration tends to 0,
given the $n_\ell$ observations (Rubin 1981).
\end{enumerate}
In mean mode, $w_i=1/n_\ell$ and the dataset uses the empirical leaf law.
\end{proposition}
\begin{proof}
(1) $P(Y=v)=\mathbb E[\sum_{i:y_i=v}w_i]=c\,\mathbb E[w_1]=c/n_\ell$, because the Dirichlet$(1,\dots,1)$ is
exchangeable with mean $1/n_\ell$. (2) $W_v=\sum_{i:y_i=v}w_i\sim\mathrm{Beta}(c,n_\ell-c)$ by the aggregation
property, which has the stated mean and variance. (3) The posterior of $\mathrm{DP}(\alpha,H)$ given the
observations is $\mathrm{DP}(\alpha H+\sum_i\delta_{y_i})$. As $\alpha\to0$ its weights on the atoms $y_i$ tend to
$\mathrm{Dirichlet}(1,\dots,1)$ (Ferguson 1973; Rubin 1981).
\end{proof}
\noindent The tree is fitted once: synthetic rows walk it with their own synthetic predictors (group, coarse group,
parent attributes and the columns before $j$), which is the sequential synthesis of Reiter (2005) and Nowok, Raab and
Dibben (2016). A row whose group the fit never saw (a sparse fact block passed as its coarse group) goes down the
larger branch at a split on the group, and the coarse-group splits then apply to it as they do to the fitted rows.
Without \texttt{smoothing}, every synthetic value of a column is an observed value of that column, as in synthpop's
default; \texttt{smoothing: density} adds Gaussian kernel noise with Silverman's bandwidth to numbers.
\setcounter{proposition}{\value{savedpropgen}}
